\documentclass[11pt,a4paper]{article}
\usepackage[T1]{fontenc}
\usepackage[utf8]{inputenc}
\usepackage[margin=2.4cm]{geometry}
\usepackage{amsmath}
\usepackage{url}

\newcommand{\code}[1]{\texttt{\small #1}}

\title{\textbf{mindlog $\cdot$ id} \\[2pt]
  \large Cryptographic Whitepaper}
\author{mindlog id Engineering \\ \small milo@mindlog.today}
\date{May 2026 -- version 1.0}

\begin{document}
\maketitle

\begin{abstract}
\noindent
\textbf{mindlog $\cdot$ id} is an open-source digital identity and end-to-end
encrypted messenger published under the AGPLv3 licence. This document specifies
the cryptographic mechanisms shipped in the current release: the
\emph{Signal-style} pairwise channel (X3DH key agreement + Double Ratchet),
the group channel (sender keys + ECDSA), the multi-device key vault
(PBKDF2 + AES-256-GCM, optionally unlocked via WebAuthn PRF), the
out-of-band identity verification ladder (safety numbers / QR), and the
ephemerality and one-shot message extensions. All primitives are NIST-standard
or IETF-standard; no custom cipher or curve is introduced. A canonical
reference implementation (\code{src/ratchet.ts}) is shared verbatim between
the web and Android clients, and is exercised by a frozen set of test
vectors (\code{test/vectors/ratchet.json}) replayed by both runtimes to
guarantee bit-for-bit interoperability.
\end{abstract}

\section{Threat model and design goals}

We assume an active network adversary with full control over the transport
layer. We additionally assume that the server (operated by the publisher or
self-hosted) is honest-but-curious: it is trusted to deliver messages in
order on a best-effort basis but is not trusted with plaintext content,
attachments, key material, or call media.

The design goals are:
\begin{itemize}
  \item \textbf{Confidentiality and authenticity} of every direct and group
    message under standard cryptographic assumptions
    (ECDLP on NIST P-256, AES-GCM security, HMAC-SHA-256 PRF).
  \item \textbf{Forward secrecy:} compromise of long-term key material must
    not retroactively expose past messages.
  \item \textbf{Post-compromise security (self-healing):} after a single
    healthy round-trip following a state compromise, the channel returns
    to a secure state.
  \item \textbf{No mandatory directory:} no phone number or email is
    required to open a channel.
  \item \textbf{Auditability and portability:} a single reference algorithm
    spec, verbatim across runtimes, validated by shared test vectors.
\end{itemize}

\section{Primitives}

All primitives are NIST or IETF specified. We use the same set on web
(WebCrypto, \code{crypto.subtle}) and on Android (\code{javax.crypto} +
\code{java.security}).

\begin{itemize}
  \item \textbf{Elliptic curve:} NIST P-256 (\code{secp256r1}) for both
    key agreement (ECDH) and signatures (ECDSA).
  \item \textbf{KDF:} HKDF-SHA-256 (RFC 5869) with explicit, domain-separated
    \texttt{info} labels.
  \item \textbf{Symmetric AEAD:} AES-256-GCM with a fresh 96-bit IV per
    message and a 128-bit authentication tag.
  \item \textbf{MAC / chain advance:} HMAC-SHA-256.
  \item \textbf{Password-based KDF:} PBKDF2-HMAC-SHA-256 with $600\,000$
    iterations (OWASP 2023+ recommendation) for the vault unlock path.
  \item \textbf{ECDSA encoding:} fixed-length P1363 ($r \| s$, 64 bytes)
    for cross-runtime interoperability.
\end{itemize}

\section{Pairwise channel}

Each pair of identities runs an instance of the Signal-family
\textbf{Double Ratchet} initialised by \textbf{X3DH}.

\subsection{Identity and prekeys}

Every account publishes:
\begin{itemize}
  \item an identity public key $IK_A$ (P-256),
  \item a medium-term signed prekey $SPK_A$ (rotated),
  \item a pool of one-time prekeys $\{OPK_A^{(i)}\}$.
\end{itemize}
The server stores only the \emph{public} bundle and consumes one-time
prekeys atomically, gated per requesting contact, so that a single OPK is
never handed out twice.

\subsection{X3DH initial key agreement}

To open a channel toward $B$, the initiator $A$ generates an ephemeral
$EK_A$ and computes:
\begin{align*}
DH_1 &= \text{ECDH}(IK_A, SPK_B), &
DH_2 &= \text{ECDH}(EK_A, IK_B), \\
DH_3 &= \text{ECDH}(EK_A, SPK_B), &
DH_4 &= \text{ECDH}(EK_A, OPK_B^{(i)})\ \text{(if available).}
\end{align*}
The initial root key is
$RK_0 = \mathrm{HKDF}(DH_1 \| DH_2 \| DH_3 \| DH_4,\ \text{info}=\text{``mindlog/x3dh''})$.

\subsection{Double Ratchet}

State per channel: a 32-byte root key $RK$, a sending chain key $CK_s$, a
receiving chain key $CK_r$, and a Diffie--Hellman ratchet pair.

\textbf{Symmetric ratchet step.} For each message sent, the chain key is
advanced and a message key is derived:
\[
mk = \mathrm{HMAC}(CK_n, \texttt{0x01}), \qquad
CK_{n+1} = \mathrm{HMAC}(CK_n, \texttt{0x02}).
\]
The message key $mk$ is then split via HKDF into an AES-256 encryption key
and a 12-byte IV. The ciphertext is
$ct = \mathrm{AES\text{-}GCM}_{mk}(\text{header}, \text{plaintext})$,
binding the header (sender DH public, counters) as Associated Data.

\textbf{DH ratchet step.} On receipt of a message advertising a new
sender DH public $DH'$, the receiver computes $\mathrm{ECDH}(\cdot,DH')$,
mixes it into $RK$, and resets both chain keys:
\[
(RK', CK_r') = \mathrm{HKDF}(RK \| \mathrm{ECDH}(sk, DH'),\ \text{info}=\text{``mindlog/dr''}).
\]
This produces \emph{forward secrecy} (past message keys are unrecoverable
after a chain advance) and \emph{post-compromise security} (one healthy
round-trip restores secrecy after a state compromise).

\subsection{Session continuity and reset recovery}

A channel is keyed by the peer handle but \emph{bound} to the peer's
cryptographic identity, by storing a format-independent fingerprint of
$IK_{\text{peer}}$ (the raw public key). Two kinds of desynchronisation are
detected and repaired client-side; neither needs server cooperation.

\textbf{Identity change.} If the stored fingerprint no longer matches the
peer's current $IK$ (account re-created, key rotated), the local session is
discarded and a fresh X3DH handshake is run --- otherwise the counterpart
would receive messages encrypted under a dead root key and read them as
``unreadable''.

\textbf{Unilateral reset at unchanged identity.} A peer may reset its
ratchet while keeping the \emph{same} $IK$ (reinstall, cleared local store,
regenerated prekeys). The identity fingerprint is then unchanged, so the
test above does not fire, yet the two sides hold incompatible states: one
keeps sending symmetric-ratchet continuations while the other re-runs X3DH,
and \emph{both} read ``unreadable''. We therefore additionally bind a
session to its \emph{establishing ephemeral} $EK$: when an initial
(bootstrap-bearing) message arrives whose $EK$ differs from the one that
opened the current session, the receiver tears down the stale session and
replays the X3DH responder to adopt the peer's new one. An idempotent
re-send of the same bootstrap (identical $EK$, before confirmation) is
ignored, so the recovery does not loop. The conversation self-heals on the
next initiating message; only continuations already sent on the dead
session stay unrecoverable, as forward secrecy requires.

\textbf{One-time prekey replenishment.} The OPK pool is refilled to its
target from the \emph{server-reported} available count, not the local pool
size --- a client whose published OPKs were drained or wiped server-side
then always republishes, instead of computing a zero deficit and leaving
the bundle permanently OPK-less.

\subsection{Wire format}

A message is packaged as
\[
\texttt{ciphertext} = \texttt{b64url}(\text{header}) \| \texttt{"."} \| \texttt{b64}(ct),
\]
with a header version flag (\code{v=1}, \code{v=2}) for forward
compatibility. No new database column is required: the server is oblivious
to the channel state and stores opaque blobs only.

\section{Group channel}

Group conversations use \textbf{sender keys}, mirroring the Signal group
construction, with the simplification that the group admin is the
creator and membership transitions are authenticated by the admin's
pairwise channel.

Each group member $M$ holds a \emph{sender state}: a 32-byte symmetric
chain key $CK^M$ and a P-256 ECDSA signing key pair
$(sk^M_\text{sig}, pk^M_\text{sig})$.

\subsection{Sender Key Distribution Message (SKDM)}

On joining or on rotation, $M$ ships its initial $CK^M$ and
$pk^M_\text{sig}$ to every other member \emph{over their pairwise
Double Ratchet channel}. The SKDM therefore inherits the confidentiality,
authenticity, forward secrecy, and post-compromise security properties of
the pairwise channel; the server only sees a regular E2EE message tagged
with the control prefix \code{skd}.

\subsection{Encrypt and decrypt}

To send to the group, $M$ advances its chain:
\[
mk = \mathrm{HKDF}(CK^M_n, \text{info}=\text{``mindlog/sk/mk''}), \quad
CK^M_{n+1} = \mathrm{HKDF}(CK^M_n, \text{info}=\text{``mindlog/sk/ck''}).
\]
The payload is
$ct = \mathrm{AES\text{-}GCM}_{mk}(\text{header}, \text{plaintext})$,
and the entire $(\text{header}, ct)$ is signed:
$\sigma = \mathrm{ECDSA}_{sk^M_\text{sig}}(\text{header} \| ct)$,
encoded in fixed-length P1363 (64 bytes). Receivers verify $\sigma$
against the sender's registered $pk^M_\text{sig}$, then decrypt.

\subsection{Rotation}

On any membership change that reduces the visibility set (member
removal), each remaining member regenerates a fresh $CK^M$ and
re-distributes it via SKDM to every other current member. This is the
standard sender-key forward-secrecy boundary on member removal.

\section{Multi-device key vault}

The user's identity private key must travel between devices without ever
being exposed to the server in cleartext. We store it as an
\emph{encrypted vault} with two unwrap paths.

\textbf{Passphrase path.} The vault key is derived as
\[
K_\text{vault} = \mathrm{PBKDF2\text{-}HMAC\text{-}SHA256}
  (\text{passphrase},\ \text{salt}_\text{vault},\ 600\,000),
\]
producing 32 bytes; the identity key is then sealed with
$\mathrm{AES\text{-}GCM}_{K_\text{vault}}$ under a fresh 12-byte IV.
The server stores the resulting ciphertext, salt, and IV only; the
passphrase never leaves the device.

\textbf{Passkey (WebAuthn PRF) path.} Where the platform exposes the
WebAuthn \code{prf} extension, the vault is unlocked by a 32-byte PRF
output bound to the user's passkey, yielding a phishing- and
replay-resistant unlock with no shared secret to remember.

The same envelope format is used on web and Android, so a vault written
on one platform unlocks bit-identically on the other.

\section{Identity verification}

To resist a server-side man-in-the-middle attack on the initial X3DH
bundle, each pair derives a \emph{safety number}:
\[
\mathrm{SN}_{A,B} = \mathrm{SHA\text{-}256}(\min(IK_A, IK_B) \| \max(IK_A, IK_B)),
\]
encoded as a stable 60-digit string and rendered as a QR code. Two users
that scan or compare each other's safety number obtain an out-of-band
authentication of their channel that does not depend on the server's
honesty.

\section{Ephemerality, attachments, view-once}

\textbf{Disappearing messages.} The header may carry an explicit TTL
$\tau$ (seconds). Both endpoints schedule a deterministic deletion at
$t_\text{recv} + \tau$; the ciphertext on the server is purged on the
same deadline.

\textbf{View-once messages.} A view-once flag instructs the recipient
to delete plaintext, ciphertext, and header from local storage as soon
as the message is decrypted and rendered, and to acknowledge a server-
side \emph{burn} for the corresponding row. The message key is also
zeroized after a single decryption.

\textbf{Attachments and voice.} Files are encrypted in chunks with a
per-attachment AES-256-GCM key freshly derived from the per-message key
material. The server stores opaque ciphertext blobs and never sees a
filename, MIME type, or plaintext byte.

\section{Directoryless contact (invite links)}

A user $A$ may mint a single-use invite token $t$ and share the URL
\url{https://id.mindlog.today/i/<t>}. The token is bound to $A$ on the
server but does not reveal $A$'s identity attributes; upon acceptance,
both sides install a mutual \emph{contact} relation and immediately run
the X3DH exchange described above. The token is destroyed after first
use. No phone number, no email, and no directory lookup are involved.

\section{Test vectors and cross-runtime parity}

The reference algorithm lives in \code{src/ratchet.ts} and is exercised
by the frozen test vectors \code{test/vectors/ratchet.json}. The Android
runtime ships unit tests that replay the same vectors against its
\code{javax.crypto} / \code{java.security} stack
(\code{GroupSenderTest}, \code{RatchetTest}). Any algorithmic divergence
between web and Android fails CI, which is what makes the multi-client
\emph{universality} of the cryptographic core a checked property rather
than a documentation claim.

\section{Multi-device messaging}

Each account may link up to five devices. Every device generates its
own ECDH P-256 identity key pair and publishes its own X3DH prekey
bundle (SPK + OPK pool) under a stable \emph{device-id} (UUID).

\subsection*{Fan-out at send time}

When device $D_A$ of user $A$ sends a message to user $B$ (devices
$D_{B_1}, D_{B_2}, \ldots$), it:
\begin{enumerate}
  \item Fetches the prekey bundles of all active devices of $B$ \emph{and}
        of $A$'s other devices from \code{GET /api/e2e/device-prekeys/:handle}.
  \item Encrypts the plaintext independently for each target device using
        that device's X3DH bundle and per-device Double Ratchet session
        (session key: \code{dev:<deviceId>}).
  \item Posts $N$ envelopes in a single request as
        \code{\{clientMsgId, envelopes[]\}}.
\end{enumerate}

The server stores one \emph{logical} row in \code{messages} and one
\code{message\_envelopes} row per target device.  On read, each device
receives only its own envelopes; the logical metadata (timestamps,
reactions, TTL, view-once flag) is shared.

\subsection*{Device enrolment and approval}

The first device registered by an account is auto-approved.
Subsequent devices enter a \emph{pending} state until an already-approved
device calls \code{POST /api/devices/:pk/approve}.  Approval requires
an active authenticated session from an approved device (checked via the
\code{x-device-id} header); the server never auto-approves on the basis
of the access key alone.

\subsection*{Security properties}

\textbf{Forward secrecy.}  Each device maintains its own Diffie-Hellman
ratchet state; compromise of one device does not retroactively expose
sessions on other devices.

\textbf{Device addition visibility.}  Registering a new device emits an
in-app notification to all approved devices of the account, so a
fraudulent enrolment is immediately visible.  The safety number (§7)
is per-identity; pairing partners are encouraged to re-verify after any
device change.

\textbf{Group sender keys.}  The SKDM (Sender Key Distribution Message)
is itself delivered via fan-out: each group member's SKDM is encrypted
for all their active devices, so every device of every member can decrypt
group messages natively.

\subsection*{What is explicitly not preserved}

A newly enrolled device cannot read messages sent before its enrolment
(no history transfer).  The per-identity ratchet cache
(\code{/api/e2e/cache}) introduced as a temporary multi-device bridge
has been removed; fan-out is the sole mechanism.

\section{What is not done}

We do not claim post-quantum security in this release. We do not claim
metadata privacy: the server learns who talks to whom and when,
although not what is said. There is no sealed-sender construction, no
mixnet routing, and no PIR-based contact discovery. These are tracked
items on the roadmap.

\section*{References}
\begin{itemize}
  \item Marlinspike, M. and Perrin, T. \emph{The X3DH Key Agreement
    Protocol.} Signal Foundation, 2016. \url{https://signal.org/docs/specifications/x3dh/}
  \item Perrin, T. and Marlinspike, M. \emph{The Double Ratchet
    Algorithm.} Signal Foundation, 2016.
    \url{https://signal.org/docs/specifications/doubleratchet/}
  \item Signal Foundation. \emph{Sender Keys for Group Messaging.}
  \item W3C. \emph{Web Authentication, Level 3 -- PRF Extension.}
    \url{https://www.w3.org/TR/webauthn-3/}
\end{itemize}

\end{document}
